\chapter{Прямоугольник с активной алгеброй поколений: точный запрет для спектрального вакуума}

Предыдущая проверка показала, что алгебра поколений должна действовать
непосредственно на ребре цикла. Здесь строится минимальный такой
прямоугольник и проверяется его обычная односледовая спектральная динамика.

\section{Минимальный частичный блок}

Рассмотрим четыре чередующиеся вершины Краевского с общей координатой
\(G=M_3(\mathbb R)\) на одном ребре. Положительное и отрицательное
частичные подпространства имеют размерность четыре. Условие первого порядка
и лемма Шура дают внедиагональную матрицу
\begin{equation}
 M=\begin{pmatrix}yI_3&w\\v^\dagger&z\end{pmatrix},
 \qquad
 D_p=\begin{pmatrix}0&M\\M^\dagger&0\end{pmatrix},
 \label{eq:v5-active-family-M}
\end{equation}
где \(y,z\) --- скаляры, а \(v,w\in\mathbb C^3\) --- триплетные
связующие поля поколений. Произвольное матричное ребро Юкавы отсутствует.

Квартичный член, чувствительный к циклу, равен
\begin{align}
 \frac12\Tr D_p^4={}&3|y|^4+|z|^4+\|v\|^4+\|w\|^4\\
 &+2(|y|^2+|z|^2)(\|v\|^2+\|w\|^2)
 +4\Re(yz\,w^\dagger v).
 \label{eq:v5-rectangle-quartic}
\end{align}
Следовательно, член относительной ориентации действительно существует.

\section{Односледовый спектральный потенциал}

Для
\begin{equation}
 V=-\mu^2\Tr D_p^2+\lambda\Tr D_p^4,
 \qquad \lambda>0,
\end{equation}
обозначим \(r^2=\mu^2/(2\lambda)\) и \(Q=MM^\dagger\). Тогда точно
\begin{equation}
 V=2\lambda\Tr(Q-r^2I_4)^2+\text{пост.}
 \label{eq:v5-rectangle-square}
\end{equation}
В глобальном минимуме требуется \(MM^\dagger=r^2I_4\).

Верхний левый блок равен \(|y|^2I_3+ww^\dagger\). Матрица ранга один
\(ww^\dagger\) не может быть ненулевым кратным \(I_3\), поэтому \(w=0\).
Внедиагональное уравнение даёт \(yv=0\). При ненулевом радиусе вакуума
\(y\ne0\), следовательно, \(v=0\), а оставшиеся условия имеют вид
\begin{equation}|y|=|z|=r.\end{equation}

\begin{theorem}[Запрет спектрального вакуума с активными поколениями]
Во всяком глобальном минимуме обычного односледового квадратично-квартичного
спектрального потенциала
\begin{equation}v=w=0.\end{equation}
Новая алгебра устраняет произвольную матричную свободу, но её собственное
каноническое спектральное действие не конденсирует триплетные связующие поля.
\end{theorem}

\section{Гессиан и плоские квартичные направления}

В нормированном вещественном вакууме
\begin{equation}
 \mu^2=2,\quad\lambda=1,\quad y=z=1,\quad v=w=0
\end{equation}
спектр вещественного гессиана по восьми переменным равен
\begin{equation}0^{\times3},\qquad16^{\times4},\qquad48.\end{equation}
Три плоских направления задаются условием \(v=-w\). Вдоль них
\begin{equation}V(t)-V(0)=4t^4>0,\end{equation}
то есть они устойчивы в четвёртом порядке около нуля, а не конденсируются.

\section{Калибровочное следствие}

Внутренние автоморфизмы \(M_3(\mathbb R)\) содержат калибровочное действие
типа \(SO(3)\) на поколениях. Поскольку \(v=w=0\), вакуум его не нарушает.
Остаются нежелательные безмассовые поколенческие калибровочные направления,
если не вводить дополнительное действие нарушения симметрии. Такое действие
и было бы независимой относительной жёсткостью, которую требовалось вывести.

\section{Связь с предыдущими запретительными результатами}

Обычные спектральные следы минимизируют положительную функцию от
\(MM^\dagger\) и предпочитают изометрический блок. Поэтому прежнее
расхождение \((\rho^2+r^2)^2\) и \((\rho^2-r^2)^2\) не было случайностью:
активный прямоугольник воспроизводит ту же структурную склонность на другом
графе. Общая теория спектрального действия изложена в
\path{arXiv:hep-th/9606001}, а ограничения конечных диаграмм --- в
\path{arXiv:hep-th/9701081}.

\section{Вывод}

\begin{protocol}
\begin{enumerate}
 \item прямоугольник с активной алгеброй поколений: \textbf{построен};
 \item скалярный коммутант поколенческого ребра: \textbf{получен};
 \item квартичный член, чувствительный к циклу: \textbf{получен};
 \item ненулевой вакуум связующих полей: \textbf{не получен};
 \item нарушение поколенческой калибровочной симметрии:
 \textbf{не получено};
 \item обычная конечная геометрия в заданной области:
 \textbf{направление закрыто};
 \item физическое замыкание: \textbf{не достигнуто}.
\end{enumerate}
\end{protocol}

Повторное открытие требует явно необычного элемента: относительной
кривизны, скручивания, вспомогательного отображения момента или нелокального
граничного функционала. Следующая проверка ---
\path{version5_nonordinary_architecture_fork_gate}.

Машинный сертификат сохранён в
\path{s2t_v5_family_algebra_rectangle_gate_results.json}.